Now I have the requirements document and rhythm principles. Let me systematically check the input text against W1–W4 criteria and produce the corrected output.

Key issues found:
1. **W1 禁止語**: 「揺らぐ」(「揺れる」系) — 1箇所
2. **W1 メタ宣言**: 「〜を確認しておこう」「〜を整理しよう」「〜を確認しよう」等 — 複数箇所
3. Minor W2/W3 issues integrated into fixes

\section{推定 : 一部から全体を, 幅をもって語る}
\label{sec:inference}

\subsection{問い : どこまで全体の話として持ち出してよいか}

100人にアンケートをとって「みんなの意見」を語る.
数日分の売上データを見て「最近のトレンド」を語る.
私たちはふだんから, 手元の一部のデータを根拠に, もっと大きな全体を推し量ろうとしている.

冷静に考えれば, これはかなり大胆な行為だ.
母集団の全員を調べたわけではない.
そのごく一部しか知らないのに, 「これでだいたいわかった」と思いたくなる.
しかし, 本当に「だいたいわかった」と言えるのだろうか.

前章では, 標本平均を繰り返し取ったときにその値がどうばらつくかを調べ, そのばらつきの大きさを標準誤差（SE）で測れることを見た.
SE を使えば「推定値 $\pm$ 約2SE」で幅を作れそうだ, という感触も得た.
しかし, その幅を実際に使って何かを主張した経験はまだない.

では, 手元の標本の数字をもとに, 全体について何が言えるのか.
一点の数字を差し出して「全体もこれくらいです」と断言してよいのか.
それとも, 何か別の語り方が要るのか.

この章を読み終えると, 次のことができるようになる.
母平均や母比率に対して\textbf{信頼区間}を構成し, その幅が何を意味するかを正確な日本語で説明できる.
「95\%信頼区間」の95\%が何の95\%かを, 他人に誤解なく伝えられる.


\subsection{素朴な方法とその限界 : 一点の数字はどこまで信用できるか}

最も素朴なやり方は, 標本の平均をそのまま母集団の平均だとみなすことだ.
手元の50件のレビュー評価の平均が4.2なら, 「全体の平均も4.2くらいだろう」と言ってしまう.
計算も簡単だし, 直感にも合う.

しかし, 前章で見たとおり, 別の50件を取り直せば平均は4.0かもしれないし, 4.4かもしれない.
標本を取り直すたびに値が変わるのだから, 1回の数字をそのまま信じるのは心もとない.
では, そのばらつきはどのくらいの大きさなのだろうか.
「50件もあれば, まあ大したことないだろう」と思うかもしれない.

\begin{ex}{10人中7人の「多数派」}{ten-of-seven}
ある施策について, 10人にアンケートをとったところ7人が「賛成」と答えた.
割合にすると70\%だ.
「過半数が賛成」と報告したくなる.

しかし, この70\%という数字がどのくらいばらつきうるかを見積もると（方法はこの章で学ぶ）, その幅はおおよそ35\%から93\%にまで広がる.
つまり, 同じ集団から別の10人を選んだら, 賛成が35\%になる可能性も排除できない.
「過半数が賛成」どころか, 反対のほうが多い結果さえありうるのだ.

10人中7人というのは, 直感では「かなりはっきりした多数派」に見える.
しかし, 10人という標本サイズでは, その直感はあてにならない.
\end{ex}

少数標本でのばらつきは, 多くの人が予想するよりはるかに大きい.
では, 数をたくさん集めれば解決するのだろうか.
たしかに標本サイズを増やせばSEは小さくなり, ばらつきは抑えられる.
しかし, 数さえ集めれば万事うまくいくわけではない.

\begin{mybox}{240万人の調査が外した話}
1936年の米国大統領選挙で, 文芸雑誌\textit{Literary Digest}は読者に調査票を送り, 約240万人の回答を集めた.
結果はランドン候補の圧勝という予測だった.
一方, 新興の世論調査会社ギャラップはわずか5万人の調査でルーズベルト候補の勝利を予測した.
ふたを開けてみれば, ルーズベルトが地滑り的に勝利した.

何が起きたのか.
\textit{Literary Digest}は電話帳と自動車登録簿から回答者を選んだ.
1936年当時, 電話や自動車を持っている層は裕福な人に偏っていた.
240万人を集めても, 母集団を正しく映していなければ意味がない.
標本サイズを増やせばSEは小さくなり, 信頼区間は狭くなる.
しかし, 標本に偏りがあると, その「狭い幅」で堂々と的を外すことになる.
\end{mybox}

ここまでで, 2つの限界が見えた.
標本が少なければ, 一点の数字は予想以上にばらつく.
標本を増やしても, 集め方に偏りがあれば的外れになる.

数を集めること自体は大事だが, それだけでは足りない.
ばらつきの大きさを見積もったうえで, そのばらつきを織り込んだ語り方が必要になる.


\subsection{理論の最小核 : 幅はどうやって作り, その幅は何を意味するか}

\subsubsection{点推定 : まず一点で語るところから}

幅のある語り方に進む前に, 最も素朴な出発点がある.
手元の標本平均 $\bar{x}$ を使って, 母平均 $\mu$ はこのくらいだろうと一つの値を示す.
この行為を\textbf{点推定}という.

\begin{defi}{点推定}{point-estimation}
母集団の特性を表す未知の値（母数）を, 標本から計算した一つの値（統計量）で推定することを\textbf{点推定}と呼ぶ.
たとえば, 標本平均 $\bar{x}$ は母平均 $\mu$ の点推定値である.
\end{defi}

点推定は推定の出発点だ.
しかし, $\bar{x}$ は標本を取り直すたびに異なる値をとる.
「標本平均は170gでした」と言われたとして, それが母平均から1g以内に収まっているのか, 10g離れている可能性があるのか, この一文だけでは判断しようがない.
点推定値は一つの数字を差し出すだけで, その数字がどのくらい信用できるかという情報を持たないからだ.

だから, 一点の数字に加えて, ばらつきを織り込んだ「幅」が欲しくなる.

\subsubsection{SDとSEの区別 : 何のばらつきを語りたいのか}

幅を作る前に, 一つ大事な区別がある.
前章で導入した標準偏差（SD）と標準誤差（SE）のちがいだ.
どちらも「ばらつきの大きさ」を測る道具だが, 何のばらつきを測っているかがまったく異なる.

身長のデータが100人分あるとしよう.
平均は170cm, SDは10cmだった.
ここで「170 $\pm$ 20cm（つまり150cmから190cm）」と書いたとする.
この幅は, 個人の身長がどのくらいばらつくかを示している.
ある人の身長が160cmかもしれないし180cmかもしれない, という話だ.

一方, SEは $10 / \sqrt{100} = 1.0$cm だから, 「170 $\pm$ 2.0cm（つまり168cmから172cm）」と書ける.
この幅は, 母平均がどのあたりにいそうかを示している.
100人の平均が170cmだったけれど, 全員を調べたら168cmかもしれないし172cmかもしれない, という話だ.

同じデータなのに幅が10倍ちがう.
前者は\textbf{個々の値のばらつき}を語り, 後者は\textbf{母平均の推定の不確かさ}を語っている.
問いがちがうから幅もちがう.

医学論文で「患者の血圧は120 $\pm$ 15mmHg」と書かれていたとき, 15がSDなのかSEなのかで意味はまったく異なる.
SDなら「患者ごとの血圧のばらつきがこのくらいある」という報告だし, SEなら「母平均の推定幅がこのくらいだ」という報告だ.
実際の論文でもこの混同は頻繁に指摘されている.
母平均の推定幅を知りたいのだから, ここから先はSEを使って幅を作る.

\subsubsection{SEから信頼区間へ : なぜ「$\pm$ 約2」なのか}

母平均 $\mu$ を知りたいが, 手元にあるのは標本平均 $\bar{x}$ だけだ.
前章で見たとおり, $\bar{x}$ は標本を取り直すたびに異なる値をとる.
つまり, $\bar{x}$ をそのまま $\mu$ だと言い切ってしまうのは, 少し心もとない.

では, $\bar{x}$ だけを頼りに, $\mu$ の居場所をどこまで絞り込めるだろうか.

もし $\bar{x}$ が毎回まったく同じ値になるなら, 話は簡単だ.
その値がそのまま $\mu$ だと言えばいい.
しかし実際には $\bar{x}$ はばらつく.
そのばらつきの幅は, 前章で導入した標準誤差 SE で測れるのだった.
ばらつきの幅がわかっているなら, 「$\bar{x}$ は $\mu$ からそう遠くないはずだ」ということを, もう少し具体的に言えそうだ.

ここで, 前章の中心極限定理が効いてくる.
標本サイズが十分大きければ, $\bar{x}$ はおよそ正規分布に従う.
正規分布の形がわかっているおかげで, 「$\bar{x}$ が $\mu$ からどのくらい離れうるか」を確率で語れるようになる.
具体的には, $\bar{x}$ が $\mu$ を中心として $\pm 1.96 \times \mathrm{SE}$ の範囲に収まる確率がおよそ95\%だ.

この関係を $\mu$ の側から書き直すと, 次の区間が得られる.

\begin{defi}{信頼区間（正規近似の場合）}{ci-normal}
標本平均 $\bar{x}$, 標準誤差 $\mathrm{SE}$, 信頼水準95\%に対し,
\[
  \bar{x} - 1.96 \times \mathrm{SE} \;\leq\; \mu \;\leq\; \bar{x} + 1.96 \times \mathrm{SE}
\]
を母平均 $\mu$ の\textbf{95\%信頼区間}と呼ぶ.
1.96は標準正規分布の上側2.5\%点である.
\end{defi}

中心にあるのは $\bar{x}$, つまり標本から得た一点の推定値だ.
そこから左右に $1.96 \times \mathrm{SE}$ だけ広げている.
$\mathrm{SE}$ が推定値のばらつきの幅を測り, 1.96がその幅を「95\%分」に調節する係数になっている.

前章で「$\pm$ 約2SE」と書いたのは, この1.96を丸めたものだった.

区間の幅は $2 \times 1.96 \times \mathrm{SE}$ で決まる.
$\mathrm{SE}$ は標本サイズ $n$ が大きくなるほど小さくなるから, データを増やせば区間は狭くなり, $\mu$ の居場所をより精密に絞り込める.
逆に言えば, 手元のデータ数で区間がどのくらい広いかを見れば, 自分の推定がどの程度の精度なのかを判断できる.

選挙速報で, 開票率がまだ5\%なのに「当選確実」が出ることがある.
あれは出口調査の標本から信頼区間を作り, 区間の下限が50\%を超えたときに「当確」と判断している.
標本サイズが十分あれば, 全体の5\%しか開いていなくても, 母比率の居場所をかなり狭く絞り込めるのだ.

\subsubsection{95\%とは何の95\%か}

ここで「95\%の確率で $\mu$ がこの区間に入る」と読みたくなるかもしれない.
自然な読み方だし, そう言いたくなる気持ちもわかる.
しかし, それは正確ではない.

$\mu$ は未知ではあるが, 動かない一つの値だ.
動くのは $\bar{x}$ のほうであり, したがって区間のほうだ.
同じ手順で標本を取り直して区間を作ることを100回繰り返したとすると, そのうちおよそ95回の区間が $\mu$ を含み, 残りの5回は外れる.

\textbf{信頼区間とは, 「この区間に $\mu$ がいる確率」ではなく, 「この作り方で区間を作ったとき, $\mu$ を捉えられる割合」を述べたものだ}.

この区別は, 1930年代にネイマン（Jerzy Neyman）が信頼区間の枠組みを提唱したときの核心でもある.
「確率」という言葉が手順の性能を指しているのか, 今回の一本の区間を指しているのかで, 意味がまったく変わる.

ここまでの手続きには, 理論が決めた部分と人間が選ぶ部分がある.
「$\bar{x}$ が正規分布に従う」という部分は, 中心極限定理から導かれる理論的な帰結だ.
ここに人間の裁量は入っていない.
一方, 「何\%にするか」は使う側が決めることだ.
95\%なら $\pm 1.96 \times \mathrm{SE}$, 99\%なら $\pm 2.58 \times \mathrm{SE}$ というように, 正規分布の形から対応する係数が決まる.
95\%を選ぶなら, 100回に5回は $\mu$ を外す区間ができることを受け入れる代わりに, 区間の幅をそこそこ狭く保てる.
99\%にすれば外れは減るが, 区間は広がる.
95\%が慣習的に多く使われるのは実用上のバランスであって, 理論的な根拠があるわけではない.

\subsubsection{$\sigma$ が未知のとき : t分布の登場}

ここまでの議論では, $\mathrm{SE} = \sigma / \sqrt{n}$ の $\sigma$（母集団の標準偏差）がわかっている前提で話を進めた.
しかし現実には, $\sigma$ を知っている場面はほとんどない.
母平均を知らないから推定しているのに, 母集団のばらつきだけは知っているというのは, 考えてみれば不自然な状況だ.

そこで $\sigma$ の代わりに標本の標準偏差 $s$ を使う.
ただし, $s$ は標本から計算した値だから, $s$ 自体にも不確実性がある.
$\sigma$ がわかっているときには存在しなかった「もう一段の不確かさ」が加わるのだ.

$\sigma$ がわかっているとき, 統計量 $(\bar{x} - \mu) / (\sigma / \sqrt{n})$ は標準正規分布に従う.
$\sigma$ を $s$ で置き換えると, この統計量は正規分布より裾が太い分布に従うようになる.
これが前章で予告した\textbf{t分布}だ.

裾が太いとはどういうことか.
極端な値が正規分布よりも出やすいということだ.
だから, 正規分布の1.96をそのまま使うと幅が狭すぎて, 95\%の的中率を維持できない.
t分布では1.96より少し大きな値を使い, その分だけ幅を広げて, $s$ で代用したぶんの不確実性を正直に反映させる.

\begin{defi}{t分布による信頼区間}{ci-t}
母平均 $\mu$ の $100(1-\alpha)$\%信頼区間を
\[
  \bar{x} \;\pm\; t_{n-1,\,\alpha/2} \;\times\; \frac{s}{\sqrt{n}}
\]
で構成する.
$t_{n-1,\,\alpha/2}$ は自由度 $n-1$ のt分布の上側 $\alpha/2$ 点であり, $s$ は標本の標準偏差である.
\end{defi}

$\bar{x}$ が中心, $s/\sqrt{n}$ がSEの推定値, $t_{n-1,\,\alpha/2}$ が幅の調節係数だ.
正規近似のときの1.96が, t値に置き換わっただけで, 骨格は同じだ.

自由度 $n-1$ は, $s$ を計算する際に平均を一つ使っているため, 自由に動ける値が $n$ 個から $n-1$ 個に減ることを反映している.

標本サイズが大きくなると, $s$ の不確実性は小さくなり, t分布は正規分布に近づく.
$n$ が30や40を超えるあたりから, 両者の差はほとんど気にならなくなる.
しかし $n$ が10や15のときには, t分布を使うことで幅が目に見えて広がる.
この広がりこそが, 手元のデータが少ないという事実を区間の幅に正直に反映した結果だ.

実際の手順は4つで済む.
推定値 $\bar{x}$ を出す.
SE $= s/\sqrt{n}$ を計算する.
t値を表やソフトから引く.
$\bar{x} \pm t \times \mathrm{SE}$ で幅を作る.
原理は奥深いが, 手順自体は一本の式で完結する.


\subsection{手法 : 平均・割合・差の区間推定}

信頼区間を構成するには, データの集め方に関する前提が必要になる.
標本が母集団から\textbf{無作為}に抽出されていること, 各観測が互いに\textbf{独立}であること, そして極端な外れ値がなく平均やSEの計算が安定していることだ.
この前提が崩れると, 数字は出せても当てにならない.
先ほどの\textit{Literary Digest}の例が, まさにそれだった.
集め方を最初に点検することが, 区間推定の第一歩になる.

\subsubsection{平均の信頼区間}

手順はすでに理論の節で見たとおりだ.
\begin{enumerate}
  \item 標本平均 $\bar{x}$ と標本標準偏差 $s$ を計算する.
  \item 標準誤差 $\mathrm{SE} = s / \sqrt{n}$ を求める.
  \item 自由度 $n-1$, 信頼水準に対応するt値を表またはソフトから得る.
  \item $\bar{x} \pm t \times \mathrm{SE}$ で区間を構成する.
\end{enumerate}

95\%信頼区間の場合, t値は自由度が大きいと1.96に近づき, 自由度が小さいと大きくなる.
たとえば自由度9で2.26, 自由度4で2.78だ.
標本が少ないほど幅が広がるのは, 少ない情報で推定しているぶん不確かさが大きいことの反映にほかならない.

\subsubsection{割合の信頼区間}

実務では「満足と答えた人の割合」や「基準を超えた製品の割合」のように, 割合を推定したい場面も多い.

標本比率を $\hat{p} = k/n$（$k$: 該当する件数, $n$: 全体の件数）とする.
割合のSEは
\[
  \mathrm{SE}(\hat{p}) \;\approx\; \sqrt{\frac{\hat{p}(1-\hat{p})}{n}}
\]
で見積もる.

この式を見てみよう.
分子の $\hat{p}(1-\hat{p})$ は, 割合が0.5のとき最大になり, 0や1に近づくほど小さくなる.
これは, 結果が半々に割れているときが最もばらつきやすく, ほぼ全員が同じ回答をしているときにはばらつきが小さくなることを表している.
分母の $n$ が大きいほどSEは小さくなる点は, 平均の場合と同じだ.

区間は $\hat{p} \pm z \times \mathrm{SE}(\hat{p})$ で作る.
$n$ が十分大きく $\hat{p}$ が端に寄りすぎていなければ, この正規近似が使える.
$\hat{p}$ が0.05を下回ったり0.95を超えたりする場合には正規近似が甘くなるが, そのときはソフトがより保守的な方法で幅を調整してくれる.

\subsubsection{2群の差の信頼区間}

「カテゴリAとカテゴリBで平均に差はあるか」という問いに答えたいとき, それぞれの群の平均の差 $\bar{x}_A - \bar{x}_B$ に対して信頼区間を作ることができる.

差のSEは
\[
  \mathrm{SE}(\bar{x}_A - \bar{x}_B) \;=\; \sqrt{\frac{s_A^2}{n_A} + \frac{s_B^2}{n_B}}
\]
で計算する.
各群のSEの二乗を足して平方根をとっている.
なぜ足すのかというと, 二つの群が独立にばらつくとき, 差のばらつきは個々のばらつきを足し合わせたもので決まるからだ.

この式は, 二つの群の分散が等しいとは仮定していない.
2群の分散が異なる状況でも安全に使えるWelch法の考え方に基づいている.
自由度の計算はやや複雑になるが, ソフトが自動で処理してくれる.

区間は
\[
  (\bar{x}_A - \bar{x}_B) \;\pm\; t \;\times\; \mathrm{SE}(\bar{x}_A - \bar{x}_B)
\]
の形になる.

ここまで, 平均, 割合, 2群の差と3つの場面を見てきた.
どの場合も, \textbf{推定値 $\pm$ t値（またはz値） $\times$ SE}という一つの型が繰り返し現れている.
推定値が何であれ, SEの計算式が変わるだけで, 区間を作る骨格は同じだ.

\subsubsection{幅を狭めたいとき}

信頼区間の幅が広すぎて, 実用的な結論を引き出しにくい場合がある.
幅を狭めるための手段は, 大きく3つある.

まず, \textbf{標本サイズ $n$ を増やす}こと.
$\mathrm{SE}$ は $1/\sqrt{n}$ に比例して縮むから, $n$ を4倍にすれば幅は半分になる.
最も確実な方法だが, コストや時間の制約がある.

次に, \textbf{ばらつき $s$ を下げる}こと.
測定方法を安定させる, 条件を揃える, 母集団を層に分けて層ごとに推定するなどの工夫で, 標本のばらつき自体を小さくできる場合がある.

最後に, \textbf{信頼水準を下げる}こと.
95\%を90\%にすれば幅は狭くなる.
ただし, 外す確率も5\%から10\%に増える.
幅の狭さと外しにくさはトレードオフの関係にあり, どちらを優先するかは分析の目的による.


\subsection{実践 : 手を動かして幅をつける}

\subsubsection{平均の信頼区間を計算する}

ある工場で生産された商品の重量を64個測定したところ, 平均170g, 標準偏差12gだった.
母平均の95\%信頼区間を求めてみよう.

まず, SE を計算する.
\[
  \mathrm{SE} = \frac{s}{\sqrt{n}} = \frac{12}{\sqrt{64}} = \frac{12}{8} = 1.5 \;\text{g}
\]

次に, t値を求める.
自由度は $n - 1 = 63$ で, 95\%信頼区間に対応するt値は約2.00だ.
自由度が60を超えると, t値は正規分布の1.96にかなり近い.

信頼区間は
\[
  170 \;\pm\; 2.00 \times 1.5 \;=\; 170 \pm 3.0
\]
となり, \textbf{167.0gから173.0g}が95\%信頼区間だ.

この数字が何を意味するか, 具体的な場面に置いてみよう.
もしこの工場の品質基準が「母平均が168gから172gの範囲に収まること」だったとする.
信頼区間は167.0gから173.0gだから, 基準の下限168gよりも下に, 上限172gよりも上に, それぞれはみ出している.
基準を満たしているとは断言しにくい.
一点の推定値170gだけを見れば基準の真ん中に収まっているのに, 幅をつけて語ると, その安心感は崩れる.

\subsubsection{信頼区間を文で読む}

この結果を人に伝えるとき, 何と言えばよいだろうか.

「母平均の95\%信頼区間は167.0gから173.0gである」がまず事実の報告だ.
この幅の意味を補足するなら, 「同じ手順で標本を取り直して区間を作ることを繰り返せば, そのうち約95\%の区間が母平均を含む. 今回の区間はその1本である」と言える.

「95\%の確率で母平均がこの中にある」とは言わない.
母平均は固定された一つの値であり, 今回の区間が当たっているか外れているかの二択だ.
95\%は, この手順の的中率を述べている.

日常会話ではそこまで厳密に言い分ける必要はないかもしれない.
しかし, 報告書や論文で信頼区間に言及するときには, この区別を意識しておくと, 読む側の誤解を防げる.

\subsubsection{割合の信頼区間を計算する}

ある商品カテゴリ80件のうち, 星評価が4.0以上のものが48件だったとする.
標本比率は $\hat{p} = 48/80 = 0.60$（60\%）だ.

SEを求める.
\[
  \mathrm{SE}(\hat{p}) = \sqrt{\frac{0.60 \times 0.40}{80}} = \sqrt{\frac{0.24}{80}} = \sqrt{0.003} \approx 0.055
\]

$n = 80$ で $\hat{p} = 0.60$ なら正規近似が使える.
$z = 1.96$ として区間を作ると,
\[
  0.60 \;\pm\; 1.96 \times 0.055 \;=\; 0.60 \pm 0.107
\]
となり, \textbf{49.3\%から70.7\%}が95\%信頼区間だ.

60\%という数字だけ見ると「かなり高い」と感じるかもしれない.
しかし幅をつけると約49\%から71\%まで広がる.
「過半数を超えているか」という問いに対しては, 区間の下限が50\%をわずかに下回っているため, 確信を持って「超えている」とは言いにくい.
冒頭の例題で見た「10人中7人」ほど極端ではないが, 80件でもまだこの程度の幅が残るのだ.

\subsubsection{2群の差の信頼区間を計算する}

カテゴリAの商品30件の平均購買数が85, 標準偏差が20, カテゴリBの商品35件の平均購買数が72, 標準偏差が18だったとする.
差の推定値は $85 - 72 = 13$ だ.

差のSEを求める.
\[
  \mathrm{SE} = \sqrt{\frac{20^2}{30} + \frac{18^2}{35}} = \sqrt{\frac{400}{30} + \frac{324}{35}} = \sqrt{13.33 + 9.26} = \sqrt{22.59} \approx 4.75
\]

Welch法による自由度はソフトに任せるが, おおよそ59程度になる.
t値は約2.00として, 信頼区間は
\[
  13 \;\pm\; 2.00 \times 4.75 \;=\; 13 \pm 9.5
\]
つまり, \textbf{3.5から22.5}が差の95\%信頼区間だ.

この区間は0をまたいでいない.
0をまたがないということは, 「カテゴリAとBの母平均に差がない」という可能性が, この幅の中に含まれていないことを意味する.
AのほうがBより購買数が多そうだ, という方向性は読み取れる.

逆に, もし区間が $-2$ から $28$ のように0をまたいでいたら, 「差がない」可能性を排除できない.
差の信頼区間が0を含むかどうかは, 次章で学ぶ検定の考え方とも深くつながっている.
ただし, 「0をまたがないから差がある」と即座に断言してよいかどうかは, もう少し慎重な議論が要る.
これは次章で正面から扱う.


\subsection{まとめと次の一手 : 幅で語った先にある問い}

\textbf{点推定}は, 標本から計算した一つの値で母数を推定する行為だ.
標本平均 $\bar{x}$ は母平均 $\mu$ の点推定値であり, 推定の出発点になる.
しかし, 一点の数字だけではその推定がどのくらい信用できるかがわからない.

\textbf{信頼区間}は, 点推定に「幅」を添えることで, 推定の不確かさを目に見える形にする.
幅の物差しはSE, 幅の作り方は $\bar{x} \pm t \times \mathrm{SE}$ だ.
95\%信頼区間の「95\%」は, この手順を繰り返したときに母数を捉えられる割合を指す.
「この区間に母数がいる確率」ではない.

平均, 割合, 2群の差のいずれの場合も, 「推定値 $\pm$ t値 $\times$ SE」という共通の骨格で区間を構成できた.

しかし, 区間推定だけでは答えにくい問いもある.
2群の差の信頼区間が0をまたがなかったとき, 「この差は偶然のばらつきでは説明しにくい」とどこまで強く言えるのか.
信頼区間は推定の不確かさを示す道具であり, 「偶然では説明しにくい」を正面から問うには, もう一つ別の枠組みが要る.

次章では, 「もし差がなかったとしたら, 手元のデータほどの差が偶然で生じる可能性はどれくらいか」という問いに答える方法を学ぶ.
幅で語る技術の次に来るのは, 偶然では説明しにくいことを数で確かめる技術だ.
